\section{Overview: one question, two routes}
\label{sec:overview}

Consider a small C program: a loop over an array with an index
computation, and a question --- can \texttt{x == 42} hold at the failure
label? We narrate the route this question takes through the platform;
\Cref{sec:calculus} then makes the definitions precise.

\paragraph{The spine}
The program first moves from C to RISC-V. The translator on this edge is
a \emph{pinned} C compiler: a digest-locked gcc with recorded flags. We
can replay it bit-for-bit, but nobody can foresee or prove its output ---
its honest grade is $\trepr$, determinism without meaning. The RISC-V
program then moves to BTOR2, a bit-vector transition-system logic that
model checkers consume, and from there across a rule-for-rule bridge to
SMT-LIB, where bit-vector solvers decide the reachability query.

Each of these moves is a \emph{pair}: translator, source interpreter,
target interpreter, and a target-to-source interpreter that carries
behaviors backwards. When the RISC-V-to-BTOR2 translator runs, the
platform can immediately check it \emph{on this very program}: interpret
the RISC-V program directly (the left edge of a square), interpret the
BTOR2 translation and carry the resulting trace back to RISC-V
observables (the other three edges), and compare, step by step, on the
declared observables --- program counter, registers, halt flag. That
check passing is what the edge's $\tchk$ grade \emph{means}. Note
carefully what it covers: every hop \emph{from the compiled binary
onward}. The C hop itself has no square --- there is no C interpreter
and no ISA-to-C carry-back --- so the verified portion of the route
begins at the binary; the pinned compiler contributes reproducibility,
re-checked separately by a corpus-level differential against an
independent C verifier (\Cref{sec:instantiation}).

\paragraph{The branch}
RISC-V reaches BTOR2 twice. One translator was written against the prose
ISA specification; a second route factors through Sail, deriving its
semantics from the official RISC-V Sail model. On the segment where they
diverge --- from RISC-V to BTOR2 --- the two routes share no translator
and no carry-back; below the reconvergence they share the hub bridge,
so it is the diverse prefix that agreement corroborates
(\S\ref{sec:branching}). The
platform runs the question down both, and the answers must agree. If they
do, two independently derived encodings of RISC-V semantics corroborate
each other --- evidence strictly stronger than either route's own grade
(\Cref{lem:agree}). If they disagree, at least one route has a defect,
and the per-hop squares localize it: which hop, which step, which
observable (\Cref{lem:disagree,cor:localization}). The same dual-route
structure exists for AArch64 (Arm manual vs.\ Arm Sail model).

\paragraph{The answer comes home}
Suppose the SMT solver answers \textsc{sat}: the label is reachable, with
a model. The model is not the answer the user asked for --- it is a
valuation of SMT variables. The route's carry-backs now run in reverse:
the SMT model becomes a BTOR2 witness, the witness a RISC-V trace, the
trace an initial memory and register state --- concretely, the inputs of
the original C program. The platform then \emph{replays} those inputs in
the route's source-side interpreter --- for this spine, the shared
RISC-V interpreter on the compiled binary, the question's observable
read at the ISA level (\S\ref{sec:composition}) --- and watches the
failure state actually happen.
At that moment the user's trust obligations collapse: even if every
translator, both routes, and the solver were wrong, the replayed run
exhibits the behavior (\Cref{thm:existential}). Witness-carrying answers
are self-certifying; only the interpreter that replays them --- itself
differentially tested against an independently generated ISA simulator ---
remains to be believed.

Had the solver answered \textsc{unsat}, no replay could vouch for it.
That is where the grades, the branch agreement, the multi-engine
corroboration, and the independently checked certificates carry the
weight (\Cref{thm:universal}) --- and where the platform spends its
assurance budget.

\paragraph{Beyond one spine}
Nothing above is specific to C or RISC-V
(\Cref{fig:registry} draws the full registry graph). The same square contract admits
eBPF, Wasm, and EVM programs into the BTOR2 hub; Python functions and
chemical reaction networks reach the SMT-LIB hub directly; a SMILES
molecular graph reaches a molecular-formula target with no solver at all
(a \emph{compile pair} --- the calculus does not care that the languages
are far from computation, only that both have defined meaning). Each
pair declares what it preserves (its projection), what it rejects (typed,
enumerable $\unsupported$), and its grade; the registry composes routes
and reports each route's composed coverage, grade, and loss. The rest of
the paper makes each of these words precise, then measures the system.
